% MechKnow Audit — standalone \input{} file
% Requires: booktabs, array, multirow, xcolor, cleveref packages in parent document

\section{MechKnow Audit: System-Level Understanding}
\label{app:mechknow}

MechKnow \citep{tower2026mechknowledge} asks whether a body of mechanistic
findings constitutes system-level understanding, assessed along three
dimensions: \emph{parcellation coherence} (is the decomposition compatible
with other work?), \emph{individual validity} (does each claimed mechanism
pass its own validation?), and \emph{collective coverage} (what fraction of
the system's behavior is accounted for?). System-level understanding requires
all three. This appendix audits the paper against each dimension.

\begin{table}[t]
\centering
\caption{MechKnow assessment. Each dimension is rated as Strong, Moderate, or
Weak, with the key evidence and principal gap identified.}
\label{tab:mechknow}
\small
\setlength{\tabcolsep}{3.5pt}
\begin{tabular}{@{}lllp{4.2cm}@{}}
\toprule
\textbf{Dimension} & \textbf{Rating} & \textbf{Key evidence} & \textbf{Principal gap} \\
\midrule
Parcellation coherence
  & Moderate
  & Compatible with DAS \citep{geiger2023finding}
  & Incommensurable with head-level circuit analysis \\
Individual validity
  & Strong
  & MechVal: Triangulated (atlas), Mech.\ Supported (pi-SAE)
  & Cross-model validation untested \\
Collective coverage
  & Weak
  & 14 operations (synthetic), 6 GPT-2 tasks
  & No denominator; synthetic coverage undefined \\
\bottomrule
\end{tabular}
\end{table}

\paragraph{Parcellation coherence: a deliberate break.}
The paper uses Grassmannian subspaces as its unit of analysis, with DAS as
the parcellation method. This is directly compatible with the causal
abstraction framework of \citet{geiger2021causal, geiger2023finding}: both
identify causal variables as low-dimensional subspaces of the residual
stream. However, this parcellation is \emph{incommensurable} with the
head-level circuit analysis dominant in the interpretability literature
\citep{wang2022interpretability, conmy2023automated}. Head-level analyses
decompose the model into discrete components (attention heads, MLP layers)
and ask which components are causally relevant. Subspace-level analyses
decompose the \emph{representation} into geometric objects and ask which
directions are causally relevant. The two parcellations cross-cut: a single
attention head's output may contain multiple causal subspaces, and a single
causal subspace may span multiple heads' outputs.

The paper implicitly argues for subspace-level parcellation over head-level
parcellation. The subspace degeneracy result (\S\ref{sec:gauge_freedom})
provides the sharpest argument: the causal variable is robust across
exponentially many subspace parameterizations, but is not localizable to
specific components. This is precisely the parcellation coherence problem
that MechKnow diagnoses---the paper proposes a new decomposition that cannot
be directly compared to existing circuit-level findings. Whether this is a
feature (subspaces are the correct unit) or a limitation (results cannot
integrate with circuit-level work) depends on which parcellation is
ultimately more productive.

The Structured pi-SAE introduces a further parcellation layer: the latent
space is split into $z_\text{causal}$ and $z_\text{nuisance}$, with $L_1$
sparsity on the causal factor. This is compatible with the sparse dictionary
learning parcellation used by \citet{bricken2023towards} and related SAE
work, providing a bridge between the subspace-level and feature-level
decompositions. However, the bridge is not developed in the paper: the
pi-SAE's sparse features are not compared to SAE features from standard
training.

\paragraph{Individual validity: strong but uneven.}
The atlas claims reach \textbf{Triangulated} validity on the MechVal
scorecard (see MechVal scorecard, Table~\ref{tab:mechval}): multiple independent evidence lines
(IIA, equivariance, $k$-sweep profiles, stochastic grokking controls)
converge on the three-class partition. Key strengths include the stochastic
class as a natural experiment (same operation, opposite outcomes rules out
operation-level confounds) and the memorization artifact analysis (high IIA
with low equivariance rules out IIA as sufficient).

The Structured pi-SAE claims reach \textbf{Mechanistically Supported}:
necessity and sufficiency are established (the $2 \times 2$ ablation shows
both structured prior and sparsity are required), and the NL-DAS comparison
demonstrates that unconstrained nonlinear methods fail where the pi-SAE
succeeds. The gap to Triangulated is the absence of independent (non-IIA)
validation of the recovered latent structure---the pi-SAE's causal variables
are validated entirely within the interchange intervention paradigm. An
independent line of evidence (e.g., probing the latent space for known
algebraic structure, or testing whether the recovered variables predict
out-of-distribution generalization) would strengthen the case.

Cross-model validation remains untested for both claims. The atlas uses a
single architecture (single-layer transformer); the GPT-2 experiments use a
single model (GPT-2-small, layer 8). Whether the three-class partition
holds for deeper transformers, different tokenizers, or alternative
architectures is unknown. This is the most significant individual validity
gap.

\paragraph{Collective coverage: the denominator problem.}
The paper covers 14 modular arithmetic operations exhaustively and 6 GPT-2
tasks. For the synthetic setting, collective coverage is strong \emph{within
the chosen domain}: all operations are fully classified, and the stochastic
class is bounded. But the domain itself is narrow---modular arithmetic over
$\Z_{113}$ with single-layer transformers---and there is no principled
denominator for ``all possible operations'' or ``all possible models.''

For GPT-2, collective coverage is limited. Six tasks (IOI, gender bias,
hypernymy, greater-than, SVA, capitals) represent a small fraction of the
model's behavioral repertoire. The paper does not claim coverage of GPT-2's
full behavior, but MechKnow requires that the fraction be quantified or at
least bounded. How many of GPT-2's causal variables are linear? How many are
nonlinear? How many are neither? These questions are not addressed.

\paragraph{Composition criterion.}
MechKnow's composition criterion asks whether the three dimensions jointly
support system-level understanding. The paper achieves strong individual
validity and moderate parcellation coherence, but weak collective coverage.
The result is a \textbf{local} rather than \textbf{global} understanding:
within the studied operations and tasks, the findings are well-validated and
mutually reinforcing. But they do not yet compose into a system-level account
of how causal geometry works across the full space of models and tasks.

The path to system-level understanding requires two advances. First,
\emph{coverage}: extending the atlas to multi-layer transformers, additional
operation families, and a systematic GPT-2 task survey would establish
whether the three-class partition is a general principle or a feature of the
specific experimental setup. Second, \emph{parcellation integration}:
connecting subspace-level findings to head-level circuit analysis (e.g.,
which heads' outputs span the causal subspace, and whether circuit-level
ablations agree with subspace-level interventions) would resolve the
incommensurability noted above.
